\section{TF--IDF calculations under the course convention}
\label{sec:ex11}
\begin{problem}
A corpus has $N=1000$ documents. The term \texttt{red} occurs in $10$ documents and appears twice in one target document; \texttt{sistema} occurs in $100$ documents and appears ten times in that target document. Use the lecture's \emph{unsmoothed, base-10} convention:
\[
\mathrm{tf}_{t,d}=\begin{cases}1+\log_{10}(\operatorname{count}(t,d)),&\operatorname{count}(t,d)>0,\\0,&\text{otherwise},\end{cases}
\quad \mathrm{idf}_t=\log_{10}\frac{N}{\mathrm{df}_t}.
\]
\begin{enumerate}[label=\textbf{(\alph*)}]
\item Calculate both IDFs.
\item Calculate TF and TF--IDF for both terms in the specified document.
\item Explain why the term occurring more often in the document may receive less TF--IDF weight.
\item What IDF does a term receive if it occurs in every document?
\end{enumerate}
\end{problem}


\begin{important}
The original guide explicitly requires the \emph{unsmoothed base-10} course definition below, \textbf{not} Scikit-Learn's default smoothed IDF. Mixing the conventions changes the answers.
\end{important}
\[
\mathrm{tf}_{t,d}=\begin{cases}
1+\log_{10}(\operatorname{count}(t,d)),&\operatorname{count}(t,d)>0,\\
0,&\operatorname{count}(t,d)=0,
\end{cases}\qquad
\mathrm{idf}_{t}=\log_{10}\!\left(\frac{N}{\mathrm{df}_{t}}\right).
\]

\subq{(a) Calculate both inverse document frequencies}
\[
\mathrm{idf}_{\rm red}=\log_{10}\!\left(\frac{1000}{10}\right)=\log_{10}(100)=\boxed{2},
\]
\[
\mathrm{idf}_{\rm sistema}=\log_{10}\!\left(\frac{1000}{100}\right)=\log_{10}(10)=\boxed{1}.
\]
The less widespread term \texttt{red} has double the IDF.

\subq{(b) Compute term frequencies and TF--IDF weights}
\textbf{Term frequencies:}
\[
\mathrm{tf}_{\rm red}=1+\log_{10}2\approx1+0.30103=1.30103,
\]
\[
\mathrm{tf}_{\rm sistema}=1+\log_{10}10=2.
\]
\textbf{Multiply by IDF:}
\[
\begin{aligned}
\mathrm{tfidf}_{\rm red}&=1.30103\times2=\boxed{2.60206},\\
\mathrm{tfidf}_{\rm sistema}&=2\times1=\boxed{2.00000}.
\end{aligned}
\]
\begin{center}\begin{tabular}{@{}lrrrr@{}}\toprule Term&Count&DF&TF&TF--IDF\\\midrule
\texttt{red}&2&10&1.30103&2.60206\\
\texttt{sistema}&10&100&2.00000&2.00000\\\bottomrule
\end{tabular}\end{center}

\subq{(c) Why can the less frequent term obtain a larger TF--IDF?}
\texttt{sistema} occurs more often in the \emph{target document}, so it has larger TF. But \texttt{red} is rarer across the \emph{corpus}, so its IDF is larger. The latter effect dominates here: $2.60206>2$. TF--IDF balances a term's local prominence against its global commonness.

\subq{(d) IDF of a term found in every document}
If $\mathrm{df}_t=N$, then $\mathrm{idf}_t=\log_{10}(N/N)=\log_{10}(1)=\boxed{0}$. Under this definition, such a ubiquitous term has zero TF--IDF weight regardless of its positive within-document count.
\begin{answer}
$\boxed{\mathrm{tfidf}(\texttt{red})\approx2.60206>\mathrm{tfidf}(\texttt{sistema})=2.0}$; a corpus-wide term has IDF zero.
\end{answer}

\subsection*{Additional numerical example: TF--IDF document search (lecture 17)}
\textbf{Lecture problem.} Find the document most relevant to \texttt{sweet love} using TF--IDF vectors and cosine similarity. The lecture-17 corpus is
\[
\begin{aligned}
d_1&:\ \texttt{sweet sweet nurse love},\\
d_2&:\ \texttt{sweet sorrow},\\
d_3&:\ \texttt{how sweet is love},\\
d_4&:\ \texttt{nurse}.
\end{aligned}
\]
Keep the lecture's ordered vocabulary
\[V=(\texttt{sweet},\texttt{nurse},\texttt{love},\texttt{how},\texttt{sorrow},\texttt{is}).\]
\textbf{Step 1 --- document frequencies and IDF.} There are $N=4$ documents. A repeated term counts only once toward document frequency $\mathrm{df}$:
\begin{center}\small
\begin{tabular}{@{}lrrl@{}}\toprule
Term & $\mathrm{df}$ & $\mathrm{idf}=\log_{10}(4/\mathrm{df})$ & Decimal\\\midrule
\texttt{sweet} & 3 & $\log_{10}(4/3)$ & $0.124939$\\
\texttt{nurse} & 2 & $\log_{10}(2)$ & $0.301030$\\
\texttt{love} & 2 & $\log_{10}(2)$ & $0.301030$\\
\texttt{how} & 1 & $\log_{10}(4)$ & $0.602060$\\
\texttt{sorrow} & 1 & $\log_{10}(4)$ & $0.602060$\\
\texttt{is} & 1 & $\log_{10}(4)$ & $0.602060$\\\bottomrule
\end{tabular}\end{center}
Define $a=\log_{10}(4/3)\approx0.124939$, $b=\log_{10}2\approx0.301030$, and $c=\log_{10}4\approx0.602060$.

\textbf{Step 2 --- construct the query vector.} Both words occur once, so $\mathrm{tf}=1+\log_{10}(1)=1$:
\[
\boxed{q=(a,0,b,0,0,0)},\quad
\norm q_2=\sqrt{a^2+b^2}\approx0.325928.
\]
\textbf{Step 3 --- construct document vectors.} Only $d_1$ repeats a term: $\mathrm{tf}(\texttt{sweet},d_1)=1+\log_{10}(2)\approx1.301030$. Set $a_2=(1+\log_{10}2)a\approx0.162549$:
\[
\begin{aligned}
d_1&=(a_2,b,b,0,0,0),\\
d_2&=(a,0,0,0,c,0),\\
d_3&=(a,0,b,c,0,c),\\
d_4&=(0,b,0,0,0,0).
\end{aligned}
\]
\textbf{Step 4 --- compute cosine similarities.} First calculate $q^\top d$, then both norms, then their quotient. For $d_1$:
\[
\begin{aligned}
q^\top d_1&=a\,a_2+b^2\approx0.110928,\\
\norm{d_1}_2&=\sqrt{a_2^2+2b^2}\approx0.455698,\\
\operatorname{sim}_{\cos}(q,d_1)&=\frac{0.110928}{0.325928(0.455698)}\approx\boxed{0.746865}.
\end{aligned}
\]
The other dot products are $q^\top d_2=a^2$, $q^\top d_3=a^2+b^2$, and $q^\top d_4=0$.
\begin{center}\small
\begin{tabular}{@{}lrrr@{}}\toprule
Document & $q^\top d$ & $\norm d_2$ & $\operatorname{sim}_{\cos}(q,d)$\\\midrule
$d_1$ & $0.110928$ & $0.455698$ & $\mathbf{0.746865}$\\
$d_2$ & $0.015610$ & $0.614887$ & $0.077889$\\
$d_3$ & $0.106229$ & $0.911691$ & $0.357498$\\
$d_4$ & $0$ & $0.301030$ & $0$\\\bottomrule
\end{tabular}\end{center}
\textbf{Step 5 --- rank the results.}
\[
\boxed{d_1\ (0.7469)>d_3\ (0.3575)>d_2\ (0.0779)>d_4\ (0)}.
\]
$d_1$ ranks first because it contains both query terms, and \texttt{sweet} occurs twice. Although $d_3$ also contains both, its additional distinctive terms \texttt{how} and \texttt{is} increase its vector norm and reduce its cosine score. \emph{All weights use the lecture's logarithmic TF and unsmoothed base-10 IDF; cosine provides its own norm normalization.}

\subsection*{Why Bag of Words is sparse (also from lecture 17)}
The lecture emphasizes that document--term vectors may be very long yet mostly zero. For example, if $p=10{,}000$ distinct vocabulary terms exist but one document contains $100$ distinct terms, the vector has at most $100$ nonzeros: at least $9{,}900/10{,}000=99\%$ of entries are zero. This is a \emph{sparse vector}. Sparse matrices store the nonzero coordinates and values, rather than allocating space for every zero.
